\section{Hamming Autocorrelation}
\label{sec:hamming}

\subsection{The Hamming Distance Between Plans}

Given two redistricting plans $\pi, \pi' : V \to [k]$ mapping census
tracts to district indices, the \emph{Hamming distance} is:
\[
  d_{\mathrm{Ham}}(\pi, \pi') = \frac{1}{|V|} \sum_{v \in V}
  \mathbf{1}[\pi(v) \neq \pi'(v)],
\]
where we canonicalize by minimizing over permutations of $[k]$ (Hungarian
algorithm) to make the metric label-invariant.

The Hamming distance ranges from 0 (identical plans) to $(k-1)/k$
(maximally different plans for $k$ equal-sized districts).
For typical \recom{} proposals that merge and resplit two adjacent
districts, $d_{\mathrm{Ham}}(\pi_t, \pi_{t+1}) \approx 1/k$ per step.

\subsection{Lag-1 Hamming Autocorrelation}

The lag-$k$ Hamming autocorrelation of the chain is:
\[
  \ham(k) = \mathrm{Corr}(d_{\mathrm{Ham}}(\pi_0, \pi_t),\,
                            d_{\mathrm{Ham}}(\pi_0, \pi_{t+k})).
\]

Intuitively, $\ham(1)$ measures how correlated consecutive plans are:
a value near 1 means the chain is moving slowly through plan space.
For \recom{} chains, typical values are $\ham(1) \approx 0.85$--$0.95$,
reflecting the local nature of the proposal distribution.

An alternative formulation---which does not require fixing a reference
plan $\pi_0$---uses the autocorrelation of a scalar statistic $f(\pi_t)$.
This is the formulation used in the $\ess$ calculation of
Section~\ref{sec:ess}.
The Hamming-based version provides a more direct geometric picture of
chain mixing.

\subsection{Hamming Distance as a Diagnostic}

A fully-mixed chain satisfies:
\[
  \mathbb{E}[d_{\mathrm{Ham}}(\pi_t, \pi_{t+k})] \approx
  \mathbb{E}[d_{\mathrm{Ham}}(\pi, \pi')],
\]
where the expectation on the right is over pairs drawn independently from
the stationary distribution.
In practice, we monitor the running average of $d_{\mathrm{Ham}}(\pi_t, \pi_{t+k})$
across the chain and check that it stabilises.

\subsection{State-Specific Hamming Profiles}

\begin{table}[h]
\centering
\caption{Lag-1 Hamming autocorrelation $\hat{\rho}_1$ by state,
         estimated from 10{,}000-step chains.
         Higher $k$ (more districts) implies more local proposals and
         higher autocorrelation.}
\label{tab:hamming}
\begin{tabular}{lcccc}
\toprule
State & $k$ & $\hat{\rho}_1$ & Steps to $\hat{\rho} < 0.5$ & $d_{\rm Ham}$/step \\
\midrule
NC & 14 & 0.87 & $\approx$ 4{,}900 & $\approx 0.14$ \\
WI &  8 & 0.85 & $\approx$ 4{,}200 & $\approx 0.25$ \\
GA & 14 & 0.88 & $\approx$ 5{,}500 & $\approx 0.14$ \\
PA & 17 & 0.89 & $\approx$ 6{,}100 & $\approx 0.12$ \\
TX & 38 & 0.92 & $\approx$ 8{,}600 & $\approx 0.05$ \\
CA & 52 & 0.94 & $\approx$11{,}400 & $\approx 0.04$ \\
\bottomrule
\end{tabular}
\end{table}

The key pattern: the per-step Hamming distance decreases as $\approx 1/k$,
because each \recom{} step disturbs a fraction $\approx 1/k$ of tracts.
The autocorrelation at lag $k_0$ for which $d_{\rm Ham}(\pi_t, \pi_{t+k_0})$
reaches 0.5 scales as $O(k)$, requiring more steps for larger states.

\subsection{Implications for Ensemble Design}

From these profiles, a practitioner can compute the minimum ensemble
length needed for a given decorrelation target.
For statutory use under the \dia{}, we recommend reporting:
(1) $\ham(1)$, (2) $\ess$ computed from $\ham$, and (3) the $\rhat$
from ten parallel chains.
All three should meet the thresholds in Section~\ref{sec:statutory}
before an ensemble is admitted as expert evidence.
